
\documentclass[border=4pt]{standalone}
\usepackage{polyglossia}
\setdefaultlanguage{russian}
\setotherlanguage{english}
\usepackage{fontspec}
\IfFontExistsTF{CMU Serif}{\setmainfont{CMU Serif}}{\setmainfont{Liberation Serif}}
\IfFontExistsTF{CMU Sans Serif}{\setsansfont{CMU Sans Serif}}{\setsansfont{Liberation Sans}}
\IfFontExistsTF{CMU Typewriter Text}{\setmonofont{CMU Typewriter Text}}{\setmonofont{DejaVu Sans Mono}}
\usepackage{amsmath,amssymb,mathtools,bm}
\usepackage{xcolor}
\definecolor{accent}{HTML}{1F4E79}
\definecolor{accent2}{HTML}{B85450}
\definecolor{shade}{HTML}{F4F1EA}
\colorlet{prosvBlue}{accent}
\colorlet{prosvRed}{accent2}
\colorlet{prosvLight}{accent!12}
\colorlet{prosvCream}{shade}
\definecolor{prosvGold}{HTML}{E8B647}
\definecolor{prosvGreen}{HTML}{6A9F58}
\colorlet{prosvGray}{black!55}
\usepackage{tikz}
\usetikzlibrary{calc, arrows.meta, decorations.pathreplacing,
                positioning, shapes.geometric, shapes.misc, fit,
                backgrounds}
\usepackage{pgfplots}
\pgfplotsset{compat=1.18}
\newcommand{\R}{\mathbb{R}}
\newcommand{\N}{\mathbb{N}}
\newcommand{\Z}{\mathbb{Z}}
\newcommand{\eps}{\varepsilon}
\newcommand{\norm}[1]{\left\|#1\right\|}
\newcommand{\abs}[1]{\left|#1\right|}
\newcommand{\NOD}{\gcd}
\newcommand{\Eul}{\varphi}
\newcommand{\Zn}[1]{\mathbb{Z}_{#1}}
\newcommand{\modn}[1]{\,(\mathrm{mod}\,#1)}
\newcommand{\term}[1]{\textcolor{accent2}{\textbf{#1}}}
\newcommand{\engterm}[1]{\textit{#1}}
\DeclareMathOperator*{\argmin}{arg\,min}
\DeclareMathOperator*{\argmax}{arg\,max}
\begin{document}

\begin{tikzpicture}[font=\small, every node/.style={align=center}]
  
  \node[draw=prosvBlue, fill=prosvLight, rounded corners=2pt, minimum width=4cm, minimum height=1cm] (f) at (0,0) {Прямое вычисление:\\ \textit{очень быстро}};
  \node[draw=prosvRed, fill=prosvCream, rounded corners=2pt, minimum width=4cm, minimum height=1cm] (fi) at (0,-1.8) {Обращение без ключа:\\ \textit{экспоненциально долго}};
  \node[draw=prosvGreen, fill=prosvLight, rounded corners=2pt, minimum width=4cm, minimum height=1cm] (fis) at (0,-3.6) {Обращение со знанием секрета $t$:\\ \textit{полиномиально быстро}};
  
  \draw[->,>=stealth, thick, prosvBlue] (-3.0,0) -- (-2.05,0);
  \node[left, font=\scriptsize, text=prosvGray] at (-3,0) {$x$};
  \draw[->,>=stealth, thick, prosvBlue] (2.05,0) -- (3,0);
  \node[right, font=\scriptsize, text=prosvGray] at (3,0) {$f(x) = y$};
  
  \draw[->,>=stealth, thick, prosvRed] (3,-1.8) -- (2.05,-1.8);
  \node[right, font=\scriptsize, text=prosvGray] at (3,-1.8) {$y$};
  \draw[->,>=stealth, thick, prosvRed] (-2.05,-1.8) -- (-3,-1.8);
  \node[left, font=\scriptsize, text=prosvGray] at (-3,-1.8) {$x$};
  
  \draw[->,>=stealth, thick, prosvGreen] (3,-3.6) -- (2.05,-3.6);
  \node[right, font=\scriptsize, text=prosvGray] at (3,-3.6) {$y + t$};
  \draw[->,>=stealth, thick, prosvGreen] (-2.05,-3.6) -- (-3,-3.6);
  \node[left, font=\scriptsize, text=prosvGray] at (-3,-3.6) {$x$};
\end{tikzpicture}
\end{document}
